\documentclass{article}
\usepackage{amsmath,url}
\title{A maximal ordered dispatch under an unavailable local read}
\author{}\date{}
\begin{document}\maketitle
\paragraph{Decision and limits.}
At the cut, assume an actual ordered queue and a consistent snapshot of
required reads. Each job has a declared set of material effects.
An earlier unfinished job must not be overtaken on a shared effect.
The decision is which jobs may be issued now without fabricating a read or
reversing the declared precedence. Issuance is not delivery. This is one
development slice, not a claim of eventual liveness through a permanent
failure or a substitute for an independently selected mission.

\paragraph{Greatest lawful subqueue.}
Use $i\to j$ when $i$ is earlier and their effect scopes intersect.
More generally one can include a declared dependency when ordering a read
or duty, even when write addresses differ. Let $B$ be jobs with an unavailable
required input. Then $j$ can be issued exactly when no path from $B$ ends
at $j$ (include paths of length zero). If a path starts at a blocked
job, issuing its endpoint forces issuance of all predecessors along the
path, a contradiction. Conversely a vertex not reachable from $B$ has no
reachable predecessor, and is ready; the entire complement is therefore
lawful. The complement is greatest by inclusion, not merely maximal.
For the ordered effect-intersection case, one left-to-right scan maintains
the union of effects touched by held jobs: hold an item when its read fails
or its scope hits that union, then add its whole scope. Induction in queue
order implements the same reachability theorem. A dependency can be
indirect: skipping a held job does not erase a later contender that shares
an effect with one of its descendants. The finite model assumes a complete
authorized effect accounting and no writer outside the snapshot.

\paragraph{Native selected-body experiment.}
The pinned SubT broker first asks for poses of every team member. One
failed read returns before the first pending packet is consumed, even when
the unreadable member is not a party to any packet. In the fixture the
adapter reads the actual pending queue, sender identities, endpoint fanout
and the same stable pose provider. It invokes the original selected
dispatcher with only the lawful set and corresponding team references.
It shares the actual simulator radio state and application objects, not
reset sender or receiver balances. The harness fixes the relevant stamps
and RF callback; it assumes no cross-component interference, generator
state or callbacks changing hidden shared effects. No conclusions about live
Gazebo middleware readiness or stock APIs for source projection follow.

In the ordinary and fanout blocked cases the measured early output consists
of the earlier ready $a$ and unrelated $e$. The middle two pending orders
remain held by the original unavailable item. After read recovery within
the relevant rate window, the application outcomes agree with a stock
global wait except for their earlier timing. Only the first of the two
consecutive writes to the common receiver wins. If the adapter erases the
receiver dimension of $c$, it issues the later $d$ early. The same receiver
balance then prevents the promised $c$ when the read recovers. This is an
isochronous two-cut observation, not a claim that accelerating earlier
traffic is harmless against every future interleaving.

The exact cutoff of the received-byte check is part of the source law.
In the pinned routine the receiver's retained usage is compared against
the \emph{current transmitter's} configured limit. With four committed
bytes on a common receiver, two further bytes are refused from a sender
limited to five and accepted from one limited to six, for three different
receiver configuration controls. This is a dependency on both the sender
choice and shared receiver ledger. It does not imply unbounded received
traffic or any vulnerability in a deployed network.

\paragraph{Comparison.}
The selector ties direct enumeration on the same graph and information.
Action-dependency-graph execution and independence methods supply the
pertinent established abstractions
(\url{https://arxiv.org/abs/2312.04190};
\url{https://ai.dmi.unibas.ch/research/reading_group/peled-cav1993.pdf}).
The comparison isolates a source-specific advantage over the stock
global gate (two actual completions before its recovery) and an exact
read/dependency contract for reaching that advantage, not a stronger
algorithm than a matched dependency scheduler. At changing delay,
noncommuting stateful RF callbacks, or hidden recipient aliasing, the
declared graph can be wrong. Effects already handed downstream need a
terminal receipt or a justified delivery horizon, not just absence from
the broker queue. The next retirement experiment targets that boundary.
\end{document}
